% Current-I definitions and statement-only cited inputs. No Paper-II GI proof or E11 routine.

% Exact current-I input: def:common-conventions
\begin{definition}[Base and integral marking]
\hypertarget{tag-01O1}{}\label{def:common-conventions}
\StatementRecord{def:common-conventions}{proved}{01O1}{}
Let $B=\mathbb P^1_t$, $p_1=0$, $p_2=1$, $p_0=\infty$,
$\Sigma=\{p_0,p_1,p_2\}$, and $B^\circ=B\setminus\Sigma$.
The full fibres of the marked family are $3S_1$, $4S_2$ and $W$;
$S_1,S_2,W$ are reduced, and $D=W+S_1+S_2$.
We orient base meridians counterclockwise in holomorphic local
coordinates and use clockwise root deck generators
$s_j(g_jz)=e^{-2\pi i/m_j}s_j(z)$, where $(m_1,m_2)=(3,4)$ and
$g_1g_2g_0=1$. The lift of a positive base loop ends at the inverse
deck translate, and its endpoint identification applies $g_j$.

The integral homology lattice and its dual basis are
\[
 \Lambda=\mathbb Z\langle\widehat\gamma,\widehat u,
                     \widehat w,\widehat\delta\rangle,
 \qquad (\gamma,u,w,\delta).
\]
The period matrix in coordinates $(\zeta_1,\zeta_2)$ is
\begin{equation}\hypertarget{tag-03CN}{}\label{wb:pshift:eq:periods}
 \Pi_c=\begin{pmatrix}6\mu&\tau&1&0\\
                    \beta_c&\mu&0&1\end{pmatrix},
 \qquad \beta_c=\beta^{\rm part}+c,
 \qquad (\beta^{\rm part}+\tau)(p_0)=0.
\end{equation}
Its rows are $\sigma_1=6\mu\gamma+\tau u+w$ and
$\sigma_2=\beta_c\gamma+\mu u+\delta$. The monodromy matrices are
\begin{equation}\hypertarget{tag-03C8}{}\label{wb:pshift:eq:finite-data}
 A_1=\begin{pmatrix}1&0&0&0\\6&0&1&0\\-6&-1&-1&0\\-2&1&0&1\end{pmatrix},
 \quad A_2=\begin{pmatrix}1&0&0&0\\0&0&-1&0\\-6&1&0&0\\3&0&1&1\end{pmatrix},
 \quad A_0=(A_1A_2)^{-1}.
\end{equation}
They act on columns in the displayed integral basis. Homology has
integral coefficients unless specified, and complex manifolds carry
their complex orientations. Reversing the two additive coordinates
acts by $J=\left(\begin{smallmatrix}0&1\\1&0\end{smallmatrix}\right)$:
the period and Hermitian matrices become $J\Pi_c$ and
$JHJ^{\mathsf t}$, respectively.
\end{definition}

% Exact current-I input: I:def:family-parameter
\begin{definition}[The parameter and the finite translations]
\hypertarget{tag-PI01}{}\label{I:def:family-parameter}
\StatementRecord{I:def:family-parameter}{proved}{PI01}{}
Put
\begin{equation}
 D_{\rm part}=\Im\beta^{\rm part}
                  -6(\Im\mu)^2/\Im\tau,\qquad
 M=\sup D_{\rm part},\qquad
 \mathcal P=\{c\in\mathbb C:\Im c<-M\},\qquad w=e^{-\pi ic}.
\label{pub:I:eq:PD1}
\end{equation}
The scalar $w$ is the family parameter; in an exterior product,
$w$ denotes the marked covector of Definition~\ref{def:common-conventions}.
The standard affine translations are
$v_1=(1,2,-4,0)^{\mathsf t}$ and
$v_2=(-1,-3,3,0)^{\mathsf t}$, and the cusp collar is untwisted.
We write $u_{\rm Engel}=e^{-2\pi ic}=w^2$ for the multiplicative
parameter with the inverse period lift at $c=0$ as its reference.
The line-bundle identification and the effect of changing that reference
are proved in Proposition~\ref{fam:src:B.1.4}.
We use $\rho=e^{\pi i/3}$; the transformation $\tau\mapsto\tau+1$
carries the convention $e^{2\pi i/3}$ to $\rho$.
\end{definition}

% Exact current-I input: inv:B1.0
\begin{proposition}[Integral changes of marking]
\hypertarget{tag-01S2}{}\label{inv:B1.0}
\StatementRecord{inv:B1.0}{proved}{01S2}{}
In Engel's common integral basis, our marked basis has column matrix
\[
 P_{\rm mark}=\begin{pmatrix}
 0&-2&-1&-2\\2&-1&0&-1\\0&2&1&3\\-1&0&0&0
 \end{pmatrix}.
\]
His base indices $(0,1,\infty)$ correspond to $(1,2,0)$ here.
The matrices satisfy $T_0P_{\rm mark}=P_{\rm mark}A_1$ and
$T_1P_{\rm mark}=P_{\rm mark}A_2$, and his distinguished circle is
$\delta_{\rm E}=P_{\rm mark}\widehat\delta$.
The positive-root affine vectors are
$a_0=-P_{\rm mark}v_1/3$ and $a_1=-P_{\rm mark}v_2/4$.
For reversed meridians the change of marking is
\[
 Q=P_{\rm mark}S,
 \quad S=\begin{pmatrix}-1&0&0&0\\-6&0&-1&0\\6&-1&0&0\\2&1&1&1\end{pmatrix},
 \quad Q=\begin{pmatrix}2&-1&0&-2\\2&-1&0&-1\\0&2&1&3\\1&0&0&0\end{pmatrix},
\]
and $T_0Q=QA_1^{-1}$, $T_1Q=QA_2^{-1}$.
\end{proposition}

% Exact current-I input: fam:src:B.1.1
\begin{proposition}[The period quotient and its marked elliptic point]\hypertarget{tag-01Q7}{}\label{fam:src:B.1.1}\label{fam:st:B.1.1}
\hypertarget{tag-01Q9}{}\label{fam:src:B.1.3}\label{fam:st:B.1.3}
\StatementRecord{fam:src:B.1.1}{proved}{01Q7}{\cref{I:weierstrass-comparison}; \cite[Remark~3.12, p.~19]{AC2026}}
\StatementRecord{fam:src:B.1.3}{proved}{01Q9}{\cref{I:weierstrass-comparison}; \cite[Remark~3.12, p.~19]{AC2026}}

Let $E_\tau=\mathbb C/(\mathbb Z+\tau\mathbb Z)$, $P=[\mu]$, and $L_\mu=\mathcal O_{E_\tau}(P-O)$. The holomorphic map
\[
 [\xi,\eta]\longmapsto[\xi,e^{2\pi i\eta}]
\]
identifies $\mathbb C^2/\Pi_c\Lambda$ with the quotient of $L_\mu^\times$ by the lift
$A(\xi,W)=(\xi+6\mu,e^{2\pi i\beta_c}W)$. The generator in \cite[Definition~2.3, p.~5]{Engel2026} is its inverse, with translation $-6P$ and multiplier $e^{-2\pi i\beta_c}$. The section $P=[\mu]$ is a generator of the Mordell--Weil group
of the stated pointed Weierstrass surface, and
\[
 \wp_\tau(\mu)=0.
\]
Under the pointed map in the proof it is one of
$P_*=(0,\sqrt2t^2(t-1))$ and $-P_*$.
\end{proposition}

% Exact current-I input: fam:src:B.1.4
\begin{proposition}[The reference linearization]\hypertarget{tag-01QA}{}\label{fam:src:B.1.4}\label{fam:st:B.1.4}
\StatementRecord{fam:src:B.1.4}{proved}{01QA}{\SeriesRef{fam:src:B.1.3}, \SeriesRef{fam:src:B.1.1}; \cite[Lemma~2.1, pp.~3--4]{Engel2026}; \cite[Corollary~2.2, p.~4]{Engel2026}}

The line-bundle comparison of the period family with Engel's family over $B^\circ$ transports its inverse period lift to a global section of $\operatorname{Hom}(t_{6P}^*M_S,M_S)$ with divisor $F_\infty$. Take this transported section for $c=0$ as the reference linearization. Then $\lambda_0=1$ exactly and $u=e^{-2\pi ic}=w^2$. For another independently fixed reference the precise constant is its scalar ratio with this section. With the cusp branch $e^{2\pi i\mu}\to-(5+2\sqrt6)$ and $f_{\rm sc}/t\to i\sqrt3/\pi$, the section is $P=(0,-\sqrt2t^2(t-1))$.
\end{proposition}

% Exact current-I input: inv:A1.1
\begin{proposition}[Local freeness and the determinant line]\hypertarget{tag-01RR}{}\label{inv:A1.1}
\StatementRecord{inv:A1.1}{proved}{01RR}{the definitions and constructions stated in the proof}

Both $\mathcal A$ and $\mathcal E$ are locally free, of ranks three and
two.  The vertical vector $\xi=\partial/\partial\zeta_2$ is nowhere
zero as a section of $\mathcal E^\vee$.  There is an exact sequence
of holomorphic vector bundles
\begin{equation}\hypertarget{tag-01SZ}{}\label{inv:eq:A-E-line}
0\longrightarrow N\longrightarrow\mathcal E
 \xrightarrow{\langle\cdot,\xi\rangle}\OO_X\longrightarrow0,
\qquad N=\det\mathcal E=\omega_f(-2S_1-3S_2).
\end{equation}
\end{proposition}
